\section{问题重述}

\subsection{计算图在多核处理器上的执行任务}
神经网络处理器的并行计算能力只有在数据及时到达、核内空间足够且任务依赖得到满足时才能发挥。矩阵与向量操作可以分配到不同核心，也可以在同一核心的不同管道上重叠执行；张量的搬入、搬出与临时驻留则占用另一组资源。神经网络加速器架构研究将计算组织和数据流组织视为密切相关的设计问题\cite{Chen2022}。本题将这种耦合具体化为给定硬件配置下的计算图切分、核心分配与子图排序问题。

题目提供 100 个计算图以及固定的核内调度、多核模拟程序。原图由张量节点与操作节点构成，边表示数据依赖；张量大小、所在存储层、操作周期及计算管道均已给定。求解者需要将非搬运操作划入若干子图，再为各核心安排子图执行序列。输出接口只有操作到子图的映射和每核子图序列，实际指令时刻、缓存换入换出以及共享带宽分配由附件机制产生\cite{Problem2026}。

切图同时改变了三个量：能够并行执行的计算工作量、跨任务或跨核心的数据传递量，以及张量在核内占用空间的时间。子图过大容易集中负载，子图过小又会增加边界搬运和同步。因而，问题的关键在于识别哪些依赖适合保留在核内，哪些计算区域值得以通信代价换取并行执行。

\subsection{三种场景的具体要求}
问题一采用场景 A：每个子图形成独立 Task，Task 之间不能直接延续核内缓存状态。即使两个子图位于同一核心，边界数据仍须经 DDR 传递；跨核依赖还具有固定等待开销。需要在这种机制下安排切图与任务顺序，尽量缩短系统完工时间并减少额外搬运。

问题二采用场景 B：同一核心的全部子图合并为一个 Task，同核生产与消费可以复用 L1/UB 中的数据。跨核依赖通过搬出、硬件同步和搬入建立联系。合并扩大了数据复用范围，也延长了部分张量的驻留，因此分核时需要兼顾通信内化、负载均衡与缓存压力。

问题三在场景 B 上增加所有核心共享的 1 MiB 只读 FIFO Cache。读取发射时查询缓存，命中读取使用独立的 L2 带宽；完成读取后按附件规则尝试插入，容量不足时按先入先出次序淘汰。需要利用合法的切图与排序形成共享读取机会，并比较同一方案在关闭和开启 L2 时的性能。

三问均要求在 1--5 核配置下汇总加速效果，保存逐图完工时间、额外搬运量及可复现方案。第三问还需给出 Cache 命中情况。问题一、二的单核参照采用整图单核方案；第三问对该固定方案补充缓存开关比较，不进行单核方案搜索。
